\chapter{Formulazione Teorica e Framework Metodologico}
\label{ch:metodologia}

\section{Modello del Sistema e Spazi di Rappresentazione}
\label{sec:modello_sistema}

La presente sezione formalizza l'impalcatura teorica e matematica del framework analitico, definendo in modo rigoroso gli spazi vettoriali di rappresentazione e le distribuzioni di probabilità necessarie per modellare i fenomeni di transizione cross-dominio nei sistemi NIDS. 
L'obiettivo è stabilire un modello astratto, indipendente dalle specifiche implementative dei singoli dataset, in grado di garantire il rigore metodologico, la comparabilità delle metriche e l'assenza di contaminazione informativa.

A tal fine, la trattazione introduce preliminarmente la formalizzazione dei domini di origine e destinazione, delineando le regole di partizionamento preventivo e le strategie di sintesi dei benchmark ibridi. 
Successivamente, viene definito l'operatore di proiezione nello spazio delle \textit{feature} condiviso, esplicitando i criteri concettuali di filtraggio degli attributi e le proprietà di invarianza algebrica che governano la pipeline di \textit{preprocessing}.

\begin{table}[H]
\centering
\caption{Sintesi della notazione matematica e dei simboli del framework.}
\label{tab:notazione_matematica}
\small
\begin{tabular}{ll}
\hline
\textbf{Simbolo} & \textbf{Descrizione e Significato Concettuale} \\ \hline
$\mathcal{D} = (\mathcal{X}, P(\mathbf{x}))$ & Dominio di traffico (Spazio delle feature e distribuzione marginale) \\
$\mathcal{X} \subseteq \mathbb{R}^d$ & Spazio delle caratteristiche $d$-dimensionale omogeneizzato \\
$\mathbf{x} = (x_1, \dots, x_d)^\top$ & Vettore delle feature del singolo flusso di rete \\
$\mathcal{Y} = \{0, 1, \dots, C\}$ & Spazio discreto delle etichette ($0$: Normal, $>0$: Minacce) \\
$\mathcal{D}_S, \mathcal{D}_T$ & Dominio Terrestre Sorgente e Domini di Destinazione (Target) \\
$\mathcal{D}_k^{(train)}, \mathcal{D}_k^{(test)}$ & Partizioni disgiunte di addestramento e test per il dominio $k$ \\
$\boldsymbol{\psi}(\cdot)$ & Operatore di proiezione e combinazione dalle feature grezze a $\mathcal{X}$ \\
$N(\cdot)$ & Operatore di cardinalità (numero di campioni estratti) \\ \hline
\end{tabular}
\end{table}

    \subsection{Formalizzazione dei Domini e Composizione Asimmetrica del Traffico}
    \label{subsec:composizione_domini}

    La formalizzazione del framework analitico richiede la definizione rigorosa degli spazi di rappresentazione vettoriale e delle distribuzioni di probabilità associate ai diversi contesti di rete considerati. 
    Si definisce un \textit{dominio} $\mathcal{D}$ come la coppia $(\mathcal{X}, P(\mathbf{x}))$, dove $\mathcal{X} \subseteq \mathbb{R}^d$ rappresenta lo spazio delle \textit{feature} $d$-dimensionale condiviso e $P(\mathbf{x})$ la distribuzione di probabilità marginale definita sul vettore delle caratteristiche $\mathbf{x} = (x_1, x_2, \dots, x_d)^\top \in \mathcal{X}$. 
    Associato a ciascun dominio, lo spazio discreto delle etichette $\mathcal{Y}$ definisce la semantica del traffico, dove $Y = 0$ indica la classe di traffico nominale (\textit{Normal}) e $Y \in \{1, \dots, C\}$ le classi relative alle differenti famiglie di minaccia.

    \paragraph{Caratterizzazione Astratta dei Domini e Tassonomia delle Classi}
    Il framework analitico modella il fenomeno di transizione cross-dominio a partire da due macro-distribuzioni di partenza distinte\footnote{I dettagli operativi relativi ai dataset concreti impiegati per istanziare tali domini, unitamente alle metriche volumetriche esatte, sono ampiamente trattati nel capitolo successivo.}:
    \begin{enumerate}
        \item \textbf{Dominio Terrestre di Riferimento ($\mathcal{D}_S$):} Rappresenta l'ambiente stazionario di origine, contenente sia istanze di traffico nominale ($Y = 0$) sia flussi malevoli ($Y \in \mathcal{Y}_S$). 
        Per eliminare le componenti ad elevata sparsità ed esigua rappresentatività statistica, si applica una procedura di rimozione delle classi minoritarie (\textit{minority removal}), garantendo una caratterizzazione stabile del baseline di superficie.
        
        \item \textbf{Domini di Destinazione ($\mathcal{D}_T$):} Rappresentano i contesti soggetti a variazioni trasmissive o canali eterogenei, composti esclusivamente da vettori di traffico malevolo ($Y \in \mathcal{Y}_T$). 
        Tali domini vengono strutturati applicando, laddove necessario, un raggruppamento delle classi minoritarie (\textit{minority merge}) per accorpare famiglie d'attacco semanticamente affini e preservare la significatività campionaria.
    \end{enumerate}

    \paragraph{Isolamento Preventivo delle Partizioni e Politiche Anti-Data Leakage}
    Per garantire il rigore metodologico ed evitare fenomeni di contaminazione informativa tra la fase di addestramento e quella di valutazione (\textit{data leakage})\footnote{L'esecuzione dello split a monte di qualsiasi fusione previene che campioni identici della classe nominale, riutilizzati per la composizione di benchmark diversi, possano figurare contemporaneamente negli insiemi di addestramento e di test di esperimenti incrociati.}, la partizione di ciascun dominio primario $\mathcal{D}_k$ nei sotto-insiemi di addestramento ($\mathcal{D}_k^{(train)}$) e di verifica ($\mathcal{D}_k^{(test)}$) viene effettuata a monte di qualsiasi operazione di aggregazione o composizione di benchmark derivati. 
    Viene adottata una suddivisione stocastica stratificata con proporzione fissa $80/20$:
    \begin{equation}
        \mathcal{D}_k = \mathcal{D}_k^{(train)} \cup \mathcal{D}_k^{(test)}, \quad \text{con } \mathcal{D}_k^{(train)} \cap \mathcal{D}_k^{(test)} = \emptyset
    \end{equation}
    garantendo la totale indipendenza strutturale dei vettori impiegati per l'inferenza.

    \paragraph{Sintesi del Benchmark Ibrido e Proporzionamento Riconfigurato}
    A causa dell'indisponibilità di campioni di traffico nominale etichettati nei contesti di destinazione, l'impianto metodologico adotta un paradigma di \textit{costruzione ibrida del benchmark}. 
    Il traffico lecito di base viene fornito in modo coerente dalla componente nominale $P_S(\mathbf{x} \mid Y = 0)$, a cui vengono sovrapposti i flussi malevoli provenienti dalle diverse distribuzioni $P_T(\mathbf{x} \mid Y \in \mathcal{Y}_T)$.

    Per rispecchiare le condizioni operative e prevenire uno sbilanciamento artificiale dei classificatori\footnote{Nei sistemi NIDS reali, il traffico lecito costituisce la stragrande maggioranza del volume di rete. Il rapporto $10:1$ emula tale asimmetria operativa mantenendo al contempo una densità di attacchi sufficiente a garantire la stabilità statistica delle metriche di rilevamento.}, ciascun benchmark derivato impone un vincolo di proporzionamento rigido $10:1$ tra la componente benevola e la componente malevola complessiva:
    \begin{equation}
        \frac{N(Y = 0)}{\sum_{y \in \mathcal{Y}_m} N(Y = y)} = 10
    \end{equation}
    dove $N(\cdot)$ indica la cardinalità dei campioni estratti e $\mathcal{Y}_m$ l'insieme delle etichette malevole incluse nel dataset specifico.

    Nel caso dei benchmark aggregati (contenenti più famiglie di attacco contemporaneamente), la composizione interna del sottoinsieme malevolo rispetta la distribuzione naturale a priori delle classi nei domini di destinazione, evitando livellamenti artificiali e preservando la fedeltà alle distribuzioni di partenza.


    %
    %
    %
    %
    % --- INSERIRE GRAFICO DRAW.IO
    %
    %
    %
    %


    \paragraph{Tassonomia Astratta delle Configurazioni Sperimentali}
    L'applicazione di tali principi formalizza la creazione di quattro categorie funzionali di benchmark:
    \begin{itemize}
        \item \textbf{Set Nativo Sorgente:} Comprendente configurazioni aggregate e configurazioni biclasse formate dall'unione della classe nominale $Y=0$ con ciascuna delle singole classi malevole del dominio sorgente.
        \item \textbf{Set Ibridi Aggregati:} Composti dalla classe nominale $Y=0$ del dominio sorgente unita alla totalità delle classi malevole dei domini di destinazione, sia considerate congiuntamente che separatamente per segmento trasmissivo.
        \item \textbf{Set Ibridi Biclasse:} Formati dalla combinazione della classe nominale $Y=0$ con le singole famiglie d'attacco isolate all'interno dei domini target.
        \item \textbf{Set di Iniezione di Riferimento:} Progettato in conformità agli standard di letteratura, integra l'intera distribuzione nativa del dominio sorgente (traffico nominale e malevolo) con un'iniezione altamente sparsa e controllata di dati malevoli provenienti dai domini di destinazione.
    \end{itemize}

    \begin{table}[H]
    \centering
    \caption{Tassonomia astratta delle quattro categorie di benchmark sperimentali.}
    \label{tab:tassonomia_benchmark}
    \small
    \renewcommand{\arraystretch}{1.35} % Migliora la spaziatura verticale tra le righe
    \begin{tabularx}{\textwidth}{@{} >{\raggedright\arraybackslash}p{3.3cm} >{\raggedright\arraybackslash}X >{\raggedright\arraybackslash}X @{}}
    \toprule
    \textbf{Categoria Benchmark} & \textbf{Classe Benevola ($Y=0$)} & \textbf{Componente Malevola ($Y > 0$)} \\ 
    \midrule
    \textbf{Set Nativo Sorgente} & 
    Baseline nativa $P_S(\mathbf{x} \mid Y=0)$ & 
    Minacce sorgente $P_S(\mathbf{x} \mid Y \in \mathcal{Y}_S)$, modellate singolarmente o in forma aggregata. \\ 

    \textbf{Set Ibridi Aggregati} & 
    Baseline nativa $P_S(\mathbf{x} \mid Y=0)$ & 
    Totalità delle minacce target $P_T(\mathbf{x} \mid Y \in \mathcal{Y}_T)$ ricombinate con vincolo di ratio $10:1$. \\ 

    \textbf{Set Ibridi Biclasse} & 
    Baseline nativa $P_S(\mathbf{x} \mid Y=0)$ & 
    Singole famiglie di attacco target $P_T(\mathbf{x} \mid Y = y_t)$ isolate con vincolo di ratio $10:1$. \\ 

    \textbf{Set di Iniezione} & 
    Distribuzione nativa completa $\mathcal{D}_S$ ($Y=0$ e $Y \in \mathcal{Y}_S$) & 
    Iniezione altamente sparsa e controllata di vettori malevoli provenienti da $\mathcal{D}_T$. \\ 
    \bottomrule
    \end{tabularx}
    \end{table}

    \subsection{Allineamento dello Spazio delle Feature e Preprocessing Astratto}
    \label{subsec:allineamento_feature}

    La disomogeneità strutturale tra le sonde di rilevamento e i formati di tracciamento adottati nei diversi ambienti di rete comporta un'asimmetria nativa negli spazi delle caratteristiche grezze. 
    Per consentire la comparabilità cross-dominio, il framework definisce una procedura astratta di riallineamento spaziale e un impianto di \textit{preprocessing} che mappa osservazioni eterogenee in uno spazio vettoriale comune, preservando la semantica trasmissiva e ottimizzando l'efficienza computazionale.

    \paragraph{Mappatura e Omogeneizzazione dello Spazio Vettoriale}
    Siano $\mathcal{X}_S^{(grezzo)} \subseteq \mathbb{R}^{d_S}$ e $\mathcal{X}_T^{(grezzo)} \subseteq \mathbb{R}^{d_T}$ gli spazi delle caratteristiche grezze associati rispettivamente al dominio sorgente e ai domini di destinazione. 
    Per superare le differenze nomenclaturali, metriche e strutturali tra le $d_S$ e $d_T$ variabili native, si definisce un operatore di trasformazione e combinazione funzionale $\boldsymbol{\psi}: \mathbb{R}^{d_k} \rightarrow \mathbb{R}^d$ (con $k \in \{S, T\}$) tale per cui\footnote{La dimensione ridotta $d$ costituisce una costante uniforme per l'intero framework. L'esplicitazione numerica di $d$ e la definizione analitica specifica delle singole componenti $\psi_j$ sono dettagliate nel capitolo successivo.}:
    \begin{equation}
        \mathbf{x} = \boldsymbol{\psi}(\mathbf{x}^{(grezzo)}) = \left( \psi_1(\mathbf{x}^{(grezzo)}), \psi_2(\mathbf{x}^{(grezzo)}), \dots, \psi_d(\mathbf{x}^{(grezzo)}) \right)^\top \in \mathcal{X} \subseteq \mathbb{R}^d
        \label{eq:feature_mapping}
    \end{equation}
    Ciascuna componente $\psi_j$ ($j = 1, \dots, d$) rappresenta una combinazione algebrica o un'estrazione semantica costruita a partire dalle variabili primarie disponibili. 
    Tale formulazione consente di derivare un set di \textit{feature} condiviso e omogeneo $\mathcal{X}$, indipendentemente dalle difformità applicative dei formati di origine.

    \paragraph{Invarianza Concettuale e Filtraggio degli Attributi}
    La definizione dello spazio proiettato $\mathcal{X}$ risponde a rigorosi vincoli di filtraggio e armonizzazione basati su tre criteri fondamentali:
    \begin{enumerate}
        \item \textbf{Rimozione degli identificatori univoci e locali:} Gli attributi legati all'identità specifica delle sessioni (es. indirizzi IP, marcatori temporali assoluti o porte di rete puntuali) vengono esclusi a priori per prevenire fenomeni di memorizzazione o correlazioni spurie (\textit{shortcut learning})\footnote{Nei contesti NIDS, lo \textit{shortcut learning} si verifica quando il classificatore apprende associazioni mnemoniche tra identificatori locali rigidi (es. un IP specifico) e la classe malevola, perdendo qualsiasi capacità di generalizzare su topologie di rete non viste.}, garantendo che i modelli apprendano esclusivamente pattern trasmissivi e comportamentali generalizzabili.
        \item \textbf{Vincoli di calcolabilità e misurabilità cross-dominio:} Le variabili per le quali non sussiste la possibilità fisica o analitica di calcolo equivalente in entrambe le distribuzioni $P_S$ e $P_T$ vengono omesse dal set condiviso. 
        Ciò garantisce che ogni dimensione $x^{(j)} \in \mathcal{X}$ possieda un significato concettuale identico e direttamente confrontabile tra i domini.
        \item \textbf{Armonizzazione delle unità di misura ed equivalenza dimensionale:} Qualora una medesima metrica sia presente o ricavabile in entrambi i domini ma espressa mediante scala o unità di misura differenti, l'operatore $\boldsymbol{\psi}$ applica le necessarie trasformazioni dimensionali per riconducerla a una grandezza standard condivisa. 
        L'uniformazione preventiva delle unità di misura garantisce che discrepanze puramente nominali nelle scale originarie non vengano erroneamente interpretate dai modelli come fenomeno di \textit{domain shift}.
    \end{enumerate}

    \paragraph{Invarianza Monotonica e Omissione della Scalatura Esplicita}
    Un aspetto centrale del design del preprocessing riguarda la gestione della scala delle variabili. 
    Sebbene le tecniche di riscalamento lineare o di standardizzazione ($Z$-score, Min-Max scaling) siano ampiamente adottate in letteratura per uniformare i domini di variabilità, il framework teorico sfrutta le proprietà algebriche delle famiglie di classificatori impiegate.

    Per la classe dei modelli basati sul partizionamento ricorsivo dello spazio dei dati\footnote{Tale famiglia include architetture ad albero di decisione singolo e relative estensioni ad \textit{ensemble}.}, la funzione di suddivisione $s(\mathbf{x})$ lungo ciascuna dimensione $j$ poggia sull'ordinamento relativo dei valori rispetto a una soglia $\theta \in \mathbb{R}$:
    \begin{equation}
        s(\mathbf{x}) = \mathbb{I}\left(x^{(j)} \le \theta\right)
    \end{equation}
    Poiché l'indicatore di partizione $\mathbb{I}(\cdot)$ è rigorosamente invariante rispetto a qualsiasi trasformazione monotonica strettamente crescente $h: \mathbb{R} \rightarrow \mathbb{R}$, si verifica che:
    \begin{equation}
        \mathbb{I}\left(x^{(j)} \le \theta\right) \iff \mathbb{I}\left(h(x^{(j)}) \le h(\theta)\right)
    \end{equation}
    Di conseguenza, la topologia dell'albero di decisione e la costruzione dei confini di separazione risultano algebricamente identiche sia operando sulle componenti nello spazio nativo $\mathbf{x}$, sia applicando una trasformazione di scala $h(\mathbf{x})$. 
    Omettendo la fase di normalizzazione esplicita, il framework preserva la geometria naturale dei dati senza alterare la capacità discriminativa dei modelli, riducendo l'overhead computazionale nella pipeline di preprocessing.

\section{Tassonomia delle Architetture di Classificazione e Valutazione Multi-Modello}
\label{sec:tassonomia_modelli}

    \subsection{Modelli di Riferimento (Baseline) e Modelli Aggregati}
    \label{subsec:modelli_aggregati}
    % Definizione formale del modello baseline (ancora del sistema) e dei classificatori addestrati su distribuzioni aggregate multi-classe.

    \subsection{Decomposizione Specialistica: Classificatori Singoli e Biclasse}
    \label{subsec:modelli_biclasse}
    % Formalizzazione della strategia di decomposizione del problema in classificatori specializzati per singole famiglie di attacco rispetto al profilo di traffico nominale.

    \subsection{Spazio delle Famiglie di Decisione Alberiformi ed Ensembling}
    \label{subsec:famiglie_classificatori}
    % Trattazione teorica, astratta e generica sulle proprietà dei modelli di decisione basati su partizionamento dello spazio (alberi di decisione, ensemble bagging e boosting con gradiente) e sulla loro tolleranza alle distorsioni di canale.

\section{Framework di Adattamento, Granularità e Stabilità Statistica}
\label{sec:framework_adattamento}

    \subsection{Incertezza Stocastica e Valutazione Multi-Seed}
    \label{subsec:stabilita_stocastica}
    % Formalizzazione della stima dell'aspettativa e della varianza prestazionale mediante iterazioni stocastiche indipendenti (invarianza rispetto ai seed di inizializzazione).

    \subsection{Granularità della Risposta: Formulazione Binaria vs. Multiclasse}
    \label{subsec:granularita_risposta}
    % Impatto della granularità delle etichette sulla rigidità degli iperpiani discriminativi under domain shift.

\section{Framework Analitico, Spiegabilità e Significatività Statistica}
\label{sec:framework_analitico}

    \subsection{Caratterizzazione Densometrica e Proiezioni Latenti}
    \label{subsec:proiezioni_densita}
    % Modellazione teorica della stima di densità del kernel (KDE) e della riduzione dimensionale ortogonale (PCA) sui dataset di test per la visualizzazione dello scostamento distributivo.

    \subsection{Spiegabilità e Attribuzione delle Feature}
    \label{subsec:spiegabilita_shap}
    % Formalizzazione dell'attribuzione dell'importanza delle feature tramite teoria dei giochi (valori di Shapley) per l'interpretazione dei confini decisionali cross-dominio.

    \subsection{Metriche Operative, Operatività a Soglia e Significatività dei Confronti}
    \label{subsec:metriche_significativita}
    % Formalizzazione di TNR, TPR, mappe di calore, curve Precision-Recall e dinamica delle probabilità al variare delle soglie di decisione.
    % Valutazione della significatività statistica del degrado e dei confronti cross-modello tramite test d'ipotesi parametrici (Welch's t-test su F1-score).